\section{Audit, integration status, and further research}
\label{sec:audit}

\subsection{Precise disposition of the source questions}

The source baseline matters because the manuscript has already absorbed
many earlier conjectures and corrections. The audit below concerns the
relevant canonical sections and the five incoming reports at the pinned
revision; it is not a claim that every line of the collective manuscript
has been independently re-proved.

\begin{longtable}{@{}>{\raggedright\arraybackslash}p{.25\textwidth}>{\raggedright\arraybackslash}p{.43\textwidth}>{\raggedright\arraybackslash}p{.26\textwidth}@{}}
\toprule Source target & Result of this continuation & Remaining scope\\\midrule
\endhead
Nested Harmonic Jets, question 7 &
Theorem~\ref{mz:thm:zero} constructs every $F_R(-a-m;a)$.
Theorems~\ref{mz:thm:primitive} and \ref{mz:thm:coefficients} make the
coordinate construction finite and explicit. &
The smallest arithmetic or differential-algebraic coordinate class is
not proved.\\[5pt]
Nested Harmonic Jets, question 4 &
Theorems~\ref{thm:doublegerm}--\ref{thm:curved} give every transverse
ordered depth-two finite part, including curvature and reparametrization.
Theorem~\ref{thm:symcoeff} gives every fully symmetrized direction. &
Individual ordered germs in higher depth and compatible renormalization
remain separate problems.\\[5pt]
Nested Harmonic Jets, question 3 &
Independent perturbation directions have a complete finite coefficient
algorithm after full labelled symmetrization. &
This does not construct the general two-colour ordered product or every
non-diagonal Bell model.\\[5pt]
Canonical Gaussian $S_6$ and current $S_8$ candidates &
Neither short arithmetic evaluation is proved here. The current
source's conjectural labels must remain. &
A relation outside the restricted fixed-weight product presentation is
needed.\\[5pt]
Inverse-colour parity &
Section~\ref{sec:gap} supplies a two-shift disk formula, its derivative
calculus, and complete complementary-diagonal reductions. &
The underlying general parity and cyclotomic symmetry principles have
prior proofs.\\[5pt]
Polynomial spectral products &
Section~\ref{sec:spectral} proves an all-order normalized-weight calculus
and local identities, including higher poles. &
The first-derivative multiplicative anomaly retains its classical
attribution.\\
\bottomrule
\end{longtable}

\paragraph{Why the Gaussian candidates have not been promoted.}
The canonical manuscript contains an exact obstruction in its specified
level-four product-row presentation: a separating functional vanishes on
the Gaussian-supported rows but not on the mixed $S_{2m}$ coordinate
\cite[\src{chapters/05-level4-rank.tex}]{repo}.
The already proved $S_4$ identity uses a larger vocabulary, so this is
a proof-system obstruction, not an arithmetic independence theorem.
The present shifted-gap reflection belongs to the established parity
mechanism; simply specializing it does not supply the missing
weight-seven or weight-nine certificate. The source also distinguishes
its current frozen $S_8$ vector from a previously rejected vector.
Neither a tiny residual interval containing zero nor equivalence of two
candidate baskets proves the vanishing of the residual.

\subsection{A confirmed external indexing correction}

There is one concrete typographical correction in the external material.
In G\"urel \cite[Theorem 1.1, printed p.~2]{gurel}, arXiv version 2 of
3 September 2025, the second exponential sum is indexed by $\ell$ but
its denominator is printed as $k$. The original PDF, not only the HTML
conversion, contains this index. The denominator should be $\ell$.

Here is a direct derivation that fixes the correction. Under the
regularizable-spectrum hypotheses of that theorem, write
$R_\ell=\operatorname{Res}_{s=\ell}\zeta_\Lambda(s)$ and
$C_\ell=\FP_{s=\ell}\zeta_\Lambda(s)$. Near $s=0$,
\[
 \frac{(s)_\ell}{\ell!}
       =\frac{s}{\ell}+\frac{H_{\ell-1}s^2}{\ell}+O(s^3),
 \qquad
 \zeta_\Lambda(\ell+s)=\frac{R_\ell}{s}+C_\ell+O(s).
\]
The coefficient of $s$ in their product is
$(C_\ell+H_{\ell-1}R_\ell)/\ell$. Expanding the shifted spectral
series and negating its derivative therefore gives the corrected form
\begin{equation}\label{eq:Gurel-correction}
 D_\Lambda(z)=\exp\left\{-\zeta_\Lambda'(0)
  -\sum_{\ell=1}^{M}
     \frac{C_\ell+H_{\ell-1}R_\ell}{\ell}z^\ell\right\}
                                                    \Delta_\Lambda(z),
\end{equation}
where $H_0=0$, and $M$ and the canonical product $\Delta_\Lambda$
have the source's meaning. The convergent terms with $\ell>M$ form
the logarithm of that canonical product. This is an indexing correction,
not a refutation of the theorem.

\subsection{Recommended clarifications for integration}

No fresh substantive error was found in the relevant current repository
proofs examined here. The following clarifications prevent errors when
the new sections are integrated.

\begin{enumerate}[leftmargin=*,itemsep=4pt]
\item \textbf{Specify the regulator before a finite part.}
Replace any informal reference to ``the finite part at $(1,1)$'' by an
explicit path or subtraction convention. Equations
\eqref{eq:direction-FP} and \eqref{eq:curved-FP} are concrete transition
formulas between these choices.
\item \textbf{Retain the labelled symmetry factor.}
The fully symmetrized depth-$d$ value has $d!$ copies on the diagonal.
The source's individual diagonal convention is recovered by dividing by
$d!$, not by suppressing repeated permutations.
\item \textbf{Distinguish fixed and moving Gamma arguments.}
The value $F_R(-a-m;a)$ holds $z$ fixed while differentiating $s$.
The path $z=-(a+m)^s$ gives zero identically. Their derivative constraints
are compatible but are different statements.
\item \textbf{Multiply through a pole before differentiating.}
The spectral jet of order $r$ can require a coefficient of order $r+1$;
with pole order $h+1$ it can require order $r+h+1$. Equations
\eqref{sp:alljets} and \eqref{sp:loghomogeneous} show the missing terms
that a premature truncation would discard.
\item \textbf{Keep ordinary sums grouped and primitives normalized.}
The terms of \eqref{eq:product-summand} are summed together.
The combined right side of \eqref{mz:Mrec} is integrated together.
The index-one primitive \eqref{gap:eq:log-primitive} subtracts a base
point before summation. These are convergence statements in the
theorems, not discretionary numerical conventions.
\item \textbf{Preserve scope and attribution.}
Do not relabel general parity, the symmetric-sum identity, or the first
multiplicative anomaly as newly proved conjectures. Do not turn the
finite coordinate ceiling into a minimality theorem. Keep the current
$S_6$ and $S_8$ status unchanged unless a separate exact certificate is
supplied.
\end{enumerate}

\subsection{Reproducible checks and their evidential role}

The accompanying verification programs implement independent
representations where practical. The Gamma-zero calculations compare
the coordinate formula with Cauchy extraction from the original
spectral product. The shifted-gap checks compare bilateral reductions
with a Mellin integral. The anomaly checks use a finite initial segment
and an independently expanded Hurwitz tail before Cauchy extraction.
Exact polynomial checks test the finite combinatorial identities and
normalizations. The complete execution parameters, counts, residuals,
and package versions are retained in the machine-readable results.

The main numerical checks include sixteen Gamma-zero evaluations
through derivative order four, seven primitive-recursion checks,
four general shifted-gap reflections and three half-diagonal cases,
four first-derivative anomaly evaluations, and a six-term second-jet
evaluation. Observed absolute residuals in these independently evaluated
special-function checks are below $10^{-39}$; the recorded results give
the individual values and precision. These are floating-point
diagnostics, not rigorous interval certificates. Exact symbolic checks
and the analytic proofs have separate roles.

The package contains modular \LaTeX{} sources, the compiled article,
verification code, recorded results, a source-provenance manifest,
and integration notes with theorem labels. It uses no external files
from the source repository to compile. The original source snapshot is
identified by hashes rather than copied wholesale into the deliverable.

\subsection{Further research questions with concrete first targets}
\label{sec:questions}

\begin{enumerate}[leftmargin=*,label=\textbf{Q\arabic*.},itemsep=8pt]
\item \textbf{Ordered depth three, with independent exponents.}
Construct an analogue of \eqref{eq:doublegerm} for
$Z_3(1+u,1+v,1+w;a)$, displaying the intersecting polar divisors and
a normally convergent holomorphic remainder. A first deliverable is
an explicit constant term along every transverse ray, with its
symmetrization checked against \eqref{eq:sym3}. The comparison should
use the existing Laurent-expansion methods of \cite{mow}.

\item \textbf{A compatible subtraction algebra.}
Determine the local counterterms that convert the directional constants
of Sections~\ref{sec:direction}--\ref{sec:symmetric} into a specified
stuffle-preserving renormalization. The depth-two discrepancy
$-\gamma_1(a)(c/d+d/c)$ is the first necessary compatibility test.
Compare the result directly with \cite{guozhang}; path independence
must be a proved property of the chosen prescription.

\item \textbf{Tangential and nonlinear approaches to nested poles.}
Theorem~\ref{thm:curved} handles transverse analytic curves. Classify
curves tangent to $u=0$ or $u+v=0$ by their vanishing orders and give
the exact number of curve coefficients required for a finite part.
Curves lying identically in a polar divisor need an independently
specified restriction or subtraction operation.

\item \textbf{Smaller arithmetic classes for Gamma-zero coordinates.}
For a specified positive rational $x$, decide whether $D(x)$ and the
first higher coordinates $L_2(x),M_{2,2}(x)$ reduce to a stated field
of Gamma, polygamma, zeta-derivative, Dirichlet-$L$-derivative, and
algebraic-argument polylogarithm values. Fix the comparison field
before searching. Proving irreducibility requires more than failure
of symbolic simplification or an integer-relation search.

\item \textbf{Several isolated zero factors.}
Develop the counterpart of Theorem~\ref{mz:thm:zero} for a
multicolour product in which two or more factors vanish at the same
parameter. Determine how the multiplicity shifts the first nonzero
normal derivative, and how mixed spectral derivatives are encoded
by multivariate Bell polynomials. A two-colour periodic assignment
is the first case connecting directly to the source's question 3.

\item \textbf{The nonsymmetric Stieltjes sums.}
Equation~\eqref{eq:Stwo} evaluates the symmetric combination
$S_{m,n}+S_{n,m}$. Determine useful exact formulas for the difference
$S_{m,n}-S_{n,m}$, beginning with $(m,n)=(0,1)$. Through
\eqref{eq:Hcoeff}, this is exactly a question about the nonsymmetric
regular double-Hurwitz jet. Identify which part reduces to ordinary
Stieltjes constants and which part needs a new normalized integral.

\item \textbf{Primitive bases for complementary-shift polylogarithms.}
For $a\ge2$, integrate \eqref{gap:eq:diagonal} against $dz/z$ in a
specified classical or multiple-polylogarithm basis. The half-shift
formula \eqref{gap:eq:primitive} fixes the first normalization.
At rational shifts, track the root-of-unity alphabet explicitly and
separate variable-argument identities from isolated constant reductions.

\item \textbf{Noninteger indices and global shift continuation.}
Construct a contour or Mellin continuation of the two-shift reflection
in noninteger exponents. An integer-index partial fraction with a
finite summation range cannot be differentiated spectrally without
such an extension. Separately continue the shifts beyond the safe
strip, recording finite terms when lattice poles cross a contour.

\Needspace{5\baselineskip}
\item \textbf{Spectral interpolation and representation components.}
Classify finite identities between order-$r$ jets sampled at rational
normalized weights. Theorem~\ref{sp:thm:polynomial} turns this into
finite polynomial interpolation with a known top-degree correction.
Beyond alternating sums, determine which permutation-representation
components kill every nonlocal term and yield factorizations in
Vandermonde or Schur polynomials.

\item \textbf{Arithmetic and nonpolynomial multiplicities.}
Extend the all-order local calculus to quasi-polynomial or periodic
weights by residue-class splitting, and then to spectra whose base
zeta has several or higher-order poles. Section~\ref{sec:spectral}
already settles the concrete polynomial times logarithmic class.
For new classes, state the continuation and normal-convergence
hypotheses that justify every extracted coefficient.

\item \textbf{An exact certificate for $S_6$, then the current $S_8$.}
Construct a weight-seven functional equation using a larger Cayley or
iterated-integral vocabulary, reduce it to the frozen $S_6$ basket,
and exhibit a finite symbolic certificate. The source's separating
functional is a useful test that a genuinely additional relation has
entered the proof. Only after that mechanism is understood should
the weight-nine candidate be attacked by the same method. More
digits alone do not resolve the structural obstruction.

\item \textbf{Formal interfaces for finite identities and analytic limits.}
Formalize the subset telescoping law, partition M\"obius inversion,
Stirling rational-polylogarithm identity, Bell extraction, and tangent
polarization as exact algebraic components. Then formalize normal
convergence and Euler--Maclaurin remainder arguments as separate
analytic lemmas. Interval implementations of the independent numerical
checks would provide certified regression tests without being confused
with the symbolic proofs.
\end{enumerate}

